\section{Gotchas}
\label{sec:gotchas}

\begin{enumerate}[leftmargin=1.4em,itemsep=0.45em]
\item \textbf{Silent fallback when nothing is declared.} A \code{Tile} kernel with no
  \code{#[tile]} compiles, passes its numeric tests, and has no \code{tile_spec()}.
  Only \code{tests/test_tile_declarations.rs}, which parses the source, catches it.
\item \textbf{Specs are never hand-authored.} A \code{KernelTileSpec} written beside a
  kernel is decoupled from it. Seven were deleted for exactly that reason. Close a gap
  in the macro, not at the lowering site.
\item \textbf{Only bound output axes seed propagation.} A value in an output tile at an
  untiled dim is discarded. For conv2d that discards \code{B}, \code{C_OUT} and
  \code{OH}, and with them every chance of resolving \code{x}'s \code{H} window
  (\S\ref{sec:conv-overestimate}).
\item \textbf{``One per program'' has two spellings.} \code{#[tile(extent = C)]} and
  \code{#[tile(block = 1, extent = C)]} generate the same decode but different specs.
  The first is unbound (resolves to full extent), the second is bound (resolves to 1).
  BN2d inference uses the first and the BN training kernels the second.
\item \textbf{Inputs without \code{#[tile]} vanish from the spec.} conv2d's weight,
  every per-channel parameter read through a raw pointer, and BN stats' running
  buffers all contribute traffic no cost model can see.
\item \textbf{The loop index and the window must agree, unchecked.}
  \code{#[tile_loop_tile(index = ..)]} restates as expressions what \code{window(..)}
  states as names. A stride in one and not the other compiles.
\item \textbf{Declared carries can be stale.} \code{batch_norm_stats_forward}'s
  \code{#[tile_carry]} names carries the generated body does not have (\S\ref{sec:bn}).
\item \textbf{The rank in the spec must equal the rank of the edge.}
  \code{resolve_inputs} rejects a mismatch outright. Use \code{span} when the body
  walks a flattened view (\S\ref{sec:span}).
\item \textbf{Flattened tiles are products, not boxes.} Product preservation makes the
  element count right and the shape wrong once a block exceeds the innermost extent.
  A graph walk must clamp and re-box (\S\ref{sec:edges}).
\item \textbf{A walked axis is not a grid axis,} and not a binding either. The macro
  had to stop treating it as one in three places (\code{untiled_dims},
  \code{blocked_positions}, \code{grid_spec}). A kernel that walks an axis
  \emph{cannot} be split along it by the scheduler. That is the same failure as
  blocking a scan's inner axis.
\item \textbf{\code{#[tile_loop_tile]} needs exactly one blocked output axis.} The
  windowed coordinate's base is \code{__tile_range}, and the multi-blocked prelude
  names its ranges per slot. A 2-D-tiled conv (\code{BLOCK_OH} $\times$
  \code{BLOCK_OW}) cannot use the generated loop yet.
\item \textbf{Stale documentation.} Two pieces are out of date:
  \begin{itemize}
  \item \code{tile_spec.rs}'s module doc still says \code{TileWindow} ``is not yet
    consumed''. \code{resolve_inputs} consumes it.
  \item \code{contributing/TiledKernelRecipe.md} still says conv2d is on raw pointers
    and that only one axis may be blocked.
  \end{itemize}
  Trust the code, and the tests in \code{tile_resolve.rs}.
\item \textbf{Anduin is scaffolding.} \code{propagate}, \code{sub_graph_tiling} and
  \code{schedule_graph} shipped in an earlier \code{TileGraph}, but the current one has \code{from_dag}
  and \code{todo!()}s. The five \code{anduin::tests} failures on trunk are those
  stubs.
\end{enumerate}
